\section{Introduction}
\label{sec:introduction}

During earthquake loading, saturated granular soils can experience a rapid rise in pore pressure that drastically reduces their shear strength. This phenomenon, termed liquefaction, has caused severe damage to infrastructure and buildings \citep{IshiharaKoga1981,SeedIdriss1967}. A large body of laboratory work has investigated liquefaction resistance using cyclic undrained triaxial tests, examining the effects of cyclic stress ratio (CSR), relative density, and confining pressure \citep{Hyodo1991,SeedLee1966,Silver1976,Toki1986}. However, vertically propagating shear waves in the ground induce gradually varying shear stress on soil elements, leading to continuous rotation of principal stress axes \citep{Arthur1980,IshiharaTowhata1983}, which cannot be fully represented in conventional triaxial loading.

Hollow cylinder apparatus (HCA) and hollow torsional shear tests address this limitation by allowing torsional loading with principal stress axis rotation. Early studies using HCA-type devices \citep[e.g.,][]{IshiharaYasuda1975,Tatsuoka1986,YamashitaToki1993} showed that liquefaction behavior can differ from triaxial-test results and depends on the loading mode and stress path. These observations motivate the use of HCA-type boundary conditions when investigating liquefaction mechanisms under more realistic cyclic loading paths.

The influence of the initial lateral-to-vertical effective stress ratio, \Kzero{}, on liquefaction resistance has been investigated extensively, yet the reported findings vary considerably depending on test conditions. Some experimental studies reported weak or negligible dependence of liquefaction strength on the initial anisotropic stress state \citep[e.g.,][]{IshiharaTakatsu1979,YamashitaToki1993}, whereas others found clear differences between isotropically consolidated (IC) and anisotropically consolidated (AC) specimens, with the trend depending on density \citep{GeorgiannouKonstadinou2014}. \citet{Vargas2020} observed that specimens under compressive consolidation exhibited higher liquefaction resistance than those consolidated isotropically. More recently, \citet{Zhao2024} conducted systematic torsional shear tests covering both compressional and extensional consolidation states and attributed the higher liquefaction resistance under compressive consolidation to smaller principal stress axis rotation during cyclic loading. While these findings establish the macroscopic trend convincingly, the particle-scale mechanisms through which stress anisotropy alters liquefaction resistance remain to be explored.

The discrete element method \citep[DEM;][]{Cundall1979} is well suited for this problem because it provides direct access to particle-scale quantities such as contact topology, force chains, and fabric anisotropy, while avoiding some specimen-to-specimen variability associated with laboratory preparation. Numerous DEM studies have examined liquefaction-related behavior in triaxial, simple-shear, or polyhedral specimen settings \citep[e.g.,][]{Huang2018,Yang2021,Jiang2021,Morimoto2021,WuYang2022,Xie2023,YangHuang2023,Zhang2023} and have identified important microscopic factors governing liquefaction resistance. Recent DEM investigations have also addressed \Kzero{} effects: \citet{YangTaiebat2024} found that both preparation protocols and \Kzero{} influence liquefaction resistance, with the difference between IC and AC states narrowing at higher densities; \citet{Otsubo2022} showed that specimens with stronger inherent anisotropy exhibit weaker stiffness in the minor principal direction and correspondingly lower liquefaction resistance. However, triaxial configurations do not incorporate principal stress axis rotation, while periodic-boundary and polyhedral specimen settings, though capable of applying general stress paths, lack direct counterparts in standard laboratory apparatus. DEM replication of HCA tests has been developed under drained and monotonic loading, addressing principal stress rotation, shear-band formation, and stress-path effects. Existing DEM-HCA models can be grouped by the radial boundary treatment: rigid cylindrical walls \citep{Shi2021,Liu2021}, stacked-wall approximations of membrane compliance \citep{LiGuoZhang2014,LiZhangGutierrez2015,LiChenGutierrez2017}, and, more recently, fully flexible bonded-sphere membranes \citep{Song2024,Song2024suffusion}. However, extension to undrained cyclic loading under HCA-type boundary conditions remains limited.

\citet{Han2024jgs} first realized undrained cyclic torsional shear in a DEM-based HCA model by introducing a combined servo mechanism that simultaneously satisfies the HCA stress boundary conditions and the constant-volume constraint. However, their formulation assumed a constant specimen height throughout, precluding axial loading and height variation, which are commonly employed loading modes in HCA tests. \citet{Ma2024} subsequently generalized this framework to accommodate a broader range of loading conditions, including drained stress paths and combined axial--torsional loading. While these formulations successfully reproduced a variety of target stress paths, three aspects of the undrained servo problem deserve closer examination.
(a)~Under undrained loading, the constant-volume constraint is a governing equation, not an auxiliary condition. It must be substituted into the kinematic system to eliminate one degree of freedom before the stress boundary conditions can form a closed system.
(b)~The servo conditions should reproduce the laboratory total-stress boundary conditions. Because the pore water pressure $u$ is a measured response rather than a prescribed input, directly targeting individual effective stresses, as in previous formulations, is only an approximation. A valid formulation must express the servo conditions in a form that does not require $u$ as an input.
(c)~Under the constant-volume constraint, the kinematic degrees of freedom (inner radius, outer radius, and height) are geometrically coupled, so movement of one wall inevitably affects the stresses on others. The servo coefficients must therefore account for this cross-coupling, rather than being derived independently from the contact stiffness at each wall.
The present study makes this closure explicit and derives closed-form servo coefficients at every timestep by analytically inverting the coupled stiffness matrix. This analytical exactness is especially important near liquefaction, where rapid loss of particle contacts causes large timestep-to-timestep stiffness changes and accumulated servo error can compromise boundary-condition control---an instability regime also noted by \citet{Ma2024}. The formulation is then applied to systematically investigate \Kzero{} effects under undrained cyclic torsional shear with direct correspondence to laboratory HCA tests.

The present work investigates the effect of initial stress anisotropy on liquefaction resistance in DEM-simulated undrained cyclic torsional shear of hollow cylindrical specimens. The investigation is organized in two stages. First, the macroscopic \Kzero{}-dependent liquefaction resistance trend is established for representative compression, isotropic, and extension states across multiple cyclic stress ratios, accompanied by energy and stiffness analyses. Second, the parameter space is expanded to two relative densities and five \Kzero{} values, and microscopic fabric analyses in the cylindrical coordinate system are used to separate the respective roles of relative density and stress anisotropy at the particle scale.

The remainder of this paper is organized as follows. Section~\ref{sec:methodology} presents the DEM-HCA model, specimen preparation and loading protocol, calibration/validation strategy, and microscopic descriptors. Section~\ref{sec:results} reports the macroscopic and microscopic responses and discusses the mechanisms governing \Kzero{}-dependent liquefaction resistance. Section~\ref{sec:conclusions} summarizes the main findings and implications.
